\scanpage{360}
% @ head of sheet
\mspage{84}
% @ top, lines 1–6
Il suffit bien sûr de prouver ``il suffit'', donc supposons que tout objet $W$-asphérique est $W'$-asphérique, de prouver que $W\subset W'$. Soit donc $f:X\unsure{\mathbin{-}|\to}Y$ dans $W$, prouvons $f\in W'$. Pour ceci, j'u\deleted{\unsure{f}}tilise la factorisation canonique $f=pi$
% @ upper middle, factorization
\diagram{\begin{tikzcd}[ampersand replacement=\&]
X \arrow[r,"i\in W_{\unsure{\omega|W|w}}"] \& \overline X \arrow[r,"p\in\mathrm{Parf}_{W_\omega}"] \& Y{.}
\end{tikzcd}}
% @ middle, lines 8–16
On a donc aussi $i\in W$, et comme $f=pi\in W$, on a $p\in W$ donc $p\in\mathrm{Parf}_{W_\omega}\cap W$. Or $D_{W_\omega}$ signifie que cet ensemble est $\subset W^{\unsure{u|n}}$, donc on a $p\in W^u$, \deleted{\unsure{et}} en particulier les fibres de $p$ sont $W$-asphériques. Elles sont donc $W'$-asphériques, et comme $p$ est $W_\omega$-parfait et a fortiori $W'$-parfait, il s'ensuit que $p\in W'$. On a aussi $i\in W_\omega\subset W'$, donc $i\in W'$, donc aussi $f=pi\in W'$, \unsure{qed.|qcd.}

% @ lower third, last seven lines
Revenons donc à la dém. de 6.1. En vertu du lemme, il suffit de prouver que si $X$ dans $\mathrm{Cat}$ est $W$-asphérique, il est $\mathbb W_\infty$-asphérique. Mais par hyp. $\mathrm{Loc}(5)$ on sait déjà que $X$ est $0$-connexe, et choisissant un pt $x$ dans $X$, qu'il soit $\mathbb W_\infty$-asphérique s'exprime par la condition
% archive sheet number: 359
